\scanpage{12}
% archive foliation: 11
une partie compatible avec $\Sigma$, commençant avec $x_0\in\overline G(j')$, et maximale i.e. $x_N\in C_0$, prouvons $x_N\in\overline G(j')$. On \deleted{le} prouve par réc. $x_i\in\overline G(j')\quad0\leq i\leq N$.
Supposons $x_i\in\overline G(j')$, \inserted{$0\leq i\leq N-1$} prouvons $x_{i+1}\in\overline G(j')$\ldots

A vrai dire, on a fait deux fois le même raisonnement. Considérons un jeu auxiliaire $\mathcal J'=\mathcal J(j')$ : deux joueurs $j'$ et $\overline{j'}=j$, avec même ens des \deleted{situations} positions $C'=C$, \inserted{$R'=R$} et
\[
C'(j')=C(j'),\quad C'(j)=C(J')\quad(J'=J-\{j'\}),
\]
\inserted{(on pose aussi $C(j)$)}
\deleted{\unsure{Pour}}
\[
G'_0(j')=G_0(j'),\quad G'_0(j)=C'_0\setminus G'_0(j')=C_0\setminus G_0(j')
\]
\inserted{(on le note aussi $\overline G_0(j')$)}
Alors $C_0=G'_0(j)\amalg\unsure{C'_0(j')|G'_0(j')}$.

Alors les stratégies du jeu $\mathcal J$ pour $J'=J\setminus\{j'\}$ sont les stratégies du jeu $\mathcal J'=\mathcal J(j')$ pour $j$, et une stratégie \inserted{$\Sigma'$} est acceptable \inserted{pour $x\in C$} contre $j'$ ssi elle est gagnante pour $\unsure{x\in C|x}$ \inserted{(pour $j$)}.
\margin{S'il n'y a pas de parties infinies\ldots}
Donc les \inserted{positions} $x\in C$ \uline{acceptables} \deleted{pour}\inserted{contre} $j'$ sont celles qui sont gagnantes pour $j$.

Comme $G_0(j')\cap G_0(j)=\varnothing$, on sait déjà qu'un $x$ ne peut être simultanément gagnant pour $j$ et $j'$ : \inserted{gagnant pour $j$ $\Longrightarrow$ non gagnant pour $j'$} le raisonnement précédent montre \deleted{que le} l'implication en sens inverse.
